\chapter{Matrix Product Operators and Density Operators}\label{ch:mpdo}

Fix a physical dimension $d$ and a bond dimension $D$.
An MPO tensor carries two physical indices, one ket index and one bra index,
together with the virtual left and right indices.

\section{MPO tensors}

\begin{definition}[MPO tensor]\label{def:mpo_tensor}
    \lean{MPOTensor}
    \leanok
    An \emph{MPO tensor} is a family of matrices
    \[
        \{M^{ij}\}_{i,j=0}^{d-1}, \qquad M^{ij} \in \MN{D},
    \]
    indexed by a ket index $i$ and a bra index $j$. Diagrammatically,
    \[
        \TNMPOCell{M}{i}{j}
    \]
    The black node denotes the tensor $M$, the horizontal legs are virtual,
    the upper leg is the ket index $i$, and the lower leg is the bra
    index $j$.
\end{definition}

\begin{definition}[Doubled-index MPS view]\label{def:mpo_to_mps}
    \lean{MPOTensor.toMPSTensor}
    \leanok
    \uses{def:mpo_tensor}
    By combining the pair $(i,j)$ into a single physical index in
    $\{0,\ldots,d^2-1\}$, one obtains an MPS tensor with physical dimension
    $d^2$ and bond dimension $D$.
\end{definition}

\begin{definition}[Word evaluation]\label{def:mpo_eval_word}
    \lean{MPOTensor.evalWord}
    \leanok
    \uses{def:mpo_tensor}
    For words
    $u = (i_1,\ldots,i_N)$ and $v = (j_1,\ldots,j_N)$ in $\{0,\ldots,d{-}1\}^N$,
    the \emph{word evaluation} is
    \[
        M^{u,v} := M^{i_1 j_1} M^{i_2 j_2} \cdots M^{i_N j_N} \in \MN{D}.
    \]
    The empty pair of words evaluates to $\Id_D$, and word pairs of different
    lengths evaluate to $0$.
\end{definition}

\begin{definition}[MPO operator family]\label{def:mpo_family}
    \lean{MPOTensor.mpo}
    \leanok
    \uses{def:mpo_eval_word}
    The operator generated by $M$ on $N$ sites is the matrix
    $\rho^{(N)}(M)$ indexed by configurations
    $\sigma,\tau \in \{0,\ldots,d{-}1\}^N$ and given by
    \[
        \rho^{(N)}(M)_{\sigma,\tau} = \tr(M^{\sigma,\tau}).
    \]
\end{definition}

Diagrammatically,
\[
    \TNMPOChain{M}{\sigma_1}{\tau_1}{\sigma_N}{\tau_N}{N}
\]
each black node denotes the same local tensor $M$, and the matrix entry
$\rho^{(N)}(M)_{\sigma,\tau}$ is obtained by closing the outer virtual
legs cyclically and taking the resulting trace.

\begin{definition}[Hermitian MPO tensor]\label{def:mpo_hermitian}
    \lean{MPOTensor.IsHermitian}
    \leanok
    \uses{def:mpo_tensor}
    An MPO tensor is \emph{Hermitian} if
    \[
        M^{ij} = (M^{ji})^\dagger
        \qquad \text{for all } i,j \in \{0,\ldots,d{-}1\}.
    \]
\end{definition}

\section{Transfer maps}

\begin{definition}[MPO transfer map]\label{def:mpo_transfer_map}
    \lean{MPOTensor.transferMap}
    \leanok
    \uses{def:mpo_tensor}
    The \emph{transfer map} associated to $M$ is the linear map
    $\E_M : \MN{D} \to \MN{D}$ defined by
    \[
        \E_M(X) = \sum_{i,j=0}^{d-1} M^{ij} X (M^{ij})^\dagger.
    \]
\end{definition}

\begin{lemma}[Identification with the doubled-index MPS transfer map]\label{lem:mpo_transfer_eq_mps}
    \lean{MPOTensor.transferMap_eq_toMPSTensor}
    \leanok
    \uses{def:mpo_to_mps, def:mpo_transfer_map, def:transfer_map}
    The transfer map of an MPO tensor agrees with the transfer map of its
    doubled-index MPS view.
\end{lemma}

\begin{proof}
    \leanok
    \uses{def:mpo_to_mps, def:mpo_transfer_map, def:transfer_map}
    After identifying the single physical index of the doubled tensor with a
    pair $(i,j)$, the single sum defining the MPS transfer map is exactly the
    double sum defining $\E_M$.
\end{proof}

\begin{theorem}[Positivity of the MPO transfer map]\label{thm:mpo_transfer_pos}
    \lean{MPOTensor.transferMap_pos}
    \leanok
    \uses{def:mpo_transfer_map}
    If $X \ge 0$, then $\E_M(X) \ge 0$.
\end{theorem}

\begin{proof}
    \leanok
    \uses{lem:mpo_transfer_eq_mps}
    By Lemma~\ref{lem:mpo_transfer_eq_mps}, $\E_M$ is the transfer map of the
    doubled-index MPS tensor. The corresponding MPS transfer map is positive,
    so $\E_M(X) \ge 0$ whenever $X \ge 0$.
\end{proof}

\section{MPDOs and LPDOs}

\begin{definition}[MPDO]\label{def:mpdo}
    \lean{MPOTensor.IsMPDO}
    \leanok
    \uses{def:mpo_family}
    An MPO tensor is an \emph{MPDO} if
    \[
        \rho^{(N)}(M) \ge 0
        \qquad \text{for every } N \in \N.
    \]
    This is the global positivity condition of \cite[\S4.1]{Cirac2017MPDO_AnnPhys}.
\end{definition}

\begin{definition}[LPDO]\label{def:lpdo}
    \lean{MPOTensor.IsLPDO}
    \leanok
    \uses{def:mpo_tensor}
    An MPO tensor is an \emph{LPDO} if there exist an ancilla dimension $d_K$
    and matrices $A^{(i,k)} \in \MN{D}$ such that
    \[
        M^{ij} = \sum_{k=0}^{d_K-1} A^{(i,k)} (A^{(j,k)})^\dagger
        \qquad \text{for all } i,j.
    \]
    This is the local purification form of \cite[\S4.3]{Cirac2017MPDO_AnnPhys}.
    Diagrammatically, the local LPDO structure is
    \[
        \TNLPDO{A}{A^\dagger}{k}{i}{j}
    \]
    where the upper node denotes $A^{(i,k)}$, the lower node denotes
    $(A^{(j,k)})^\dagger$, the ancilla leg $k$ is contracted vertically,
    and the left and right outer legs are the MPO virtual legs.
\end{definition}

\begin{theorem}[LPDOs are Hermitian]\label{thm:lpdo_hermitian}
    \lean{MPOTensor.IsLPDO.isHermitian}
    \leanok
    \uses{def:lpdo, def:mpo_hermitian}
    Every LPDO tensor is Hermitian.
\end{theorem}

\begin{proof}
    \leanok
    \uses{def:lpdo}
    Write
    $M^{ij} = \sum_k A^{(i,k)} (A^{(j,k)})^\dagger$.
    Then
    \[
        (M^{ji})^\dagger
        = \sum_k \left(A^{(j,k)} (A^{(i,k)})^\dagger\right)^\dagger
        = \sum_k A^{(i,k)} (A^{(j,k)})^\dagger
        = M^{ij}.
    \]
\end{proof}

\begin{theorem}[LPDOs are MPDOs]\label{thm:lpdo_is_mpdo}\notready
    \lean{MPOTensor.IsLPDO.isMPDO}
    \uses{def:lpdo, def:mpdo}
    Every LPDO tensor is an MPDO.
\end{theorem}
